\section{Joint integer and residual limit theorems}\label{sec:joint}

Under the decomposition $Y_i=J_i+Z_i$ of Lemma~\ref{lem:first-alias}, write
$J_i=K_i-\E K_i$, $K=\sum_{i=1}^N K_i$, $\mu=\E K$, $T=\sum_{i=1}^N Z_i$,
$A=\Var K$, $Q=\Var T=O(1)$, $S=K-\mu+T$, and $V=\Var S$. Because the core
$I$ of size $M$ converges uniformly in $W_3$ to $\Bern(p)$ ($p=\pe$) and
carries all but $o(M)$ of the integer variance and third absolute moments,
while $\Cov(K,T)=o(\sqrt M)$, we have
\begin{equation}\label{eq:joint-moments}
 A\sim p(1-p)M,\qquad V=A+o(\sqrt A),\qquad
 \gamma=\frac{\sum_{i=1}^N\E J_i^3}{A}\longrightarrow1-2p.
\end{equation}

\begin{lemma}[Global Edgeworth and local estimates for the integer sum]\label{lem:integer}
Uniformly over all $j\in\Z$, with $x_j=(j-\mu)/\sqrt A$,
\begin{align}
 \Pp(K\le j)&=\Phi(x_j)+\frac{\phi(x_j)}{\sqrt A}
             \left\{\frac12+\frac\gamma6(1-x_j^2)\right\}
                        +o(A^{-1/2}),\label{eq:integer-edgeworth}\\
 \sqrt A\,\Pp(K=j)&=\phi(x_j)+o(1),\label{eq:integer-local}
\end{align}
and uniformly over all $k\in\Z$,
\begin{equation}\label{eq:all-layer}
                  \E\bigl[T^2\ind_{\{K=k\}}\bigr]\le \frac{K_0 Q}{\sqrt A}.
\end{equation}
All three estimates remain uniform after deleting any fixed number of
summands.
\end{lemma}
\begin{proof}
On the core $I$, every $K_i$ has atom probabilities at $0$ and $1$ bounded
below, so after deleting a bounded number of summands the remaining product
satisfies the torus damping bound
\begin{equation}\label{eq:torus-damping}
        |\chi_{\rm deleted}(u)|\le K\exp\bigl(-cM\sin^2(u/2)\bigr)
                                                   \qquad(|u|\le\pi).
\end{equation}
Expanding $\log\chi_I(u)=-Wu^2/2-\ii\kappa_Iu^3/6+O(M(\epsilon_n|u|^3+|u|^4))$
near $u=0$, inverting against the Gaussian--Hermite lattice sequence via
Euler--Maclaurin summation, and convolving with the outside block (which has
variance and third-moment sum $o(A)$) yields \eqref{eq:integer-edgeworth}
and \eqref{eq:integer-local}. For \eqref{eq:all-layer}, expanding the Fourier
transform of $T^2\ind_{\{K=k\}}$ into diagonal terms (bounded in sum by $Q$)
and off-diagonal terms (with factors $|\E[Z_ie^{\ii uJ_i}]|\le|u|\sqrt{\Var(K_i)\E Z_i^2}$)
and integrating $KQ(1+Au^2)e^{-cAu^2}$ over $[-\pi,\pi]$ gives the uniform
bound $K_0Q/\sqrt A$.
\end{proof}

\begin{lemma}[Asymptotic independence and joint local limit]\label{lem:joint-local}
Along any subsequence for which $T\Rightarrow U$, we have
$((K-\mu)/\sqrt A,T)\Rightarrow(G,U)$, where $G\sim N(0,1)$ is independent
of $U$ and $\E U=0$. Moreover, for every integer sequence $k$ with
$(k-\mu)/\sqrt A\to x\in\R$,
\begin{equation}\label{eq:joint-local}
        \sqrt A\,\Pp(K=k,\ T\in du)\Longrightarrow\phi(x)\Law(U)(du).
\end{equation}
\end{lemma}
\begin{proof}
For fixed $u,v\in\R$ and $s=u/\sqrt A$, expanding
$\sum_{i\in I}\log(\E e^{\ii sJ_i+\ii vZ_i}/\E e^{\ii vZ_i})$ through second
order in $s$ and using $\Cov(K,T)=o(\sqrt A)$ gives
$\E e^{\ii s(K-\mu)+\ii vT}=e^{-u^2/2+o(1)}\E e^{\ii vT}$, proving joint
weak convergence. To prove \eqref{eq:joint-local}, invert the Fourier
transform in the first coordinate at $k$. Since $\sup_{i\in I}\E Z_i^2\to0$,
on a fixed small cell $|s|\le c_0$ the twisted core factors
$\E e^{\ii sJ_i+\ii vZ_i}$ converge uniformly in $C^2$ to the centered
$\Bern(p)$ factor, are bounded away from zero, and have log-modulus second
derivatives at most $-c$ while their log-modulus first derivatives at $s=0$
sum to at most $K_v\sqrt A$ by Cauchy--Schwarz. Thus
\[
 \left|\prod_{i\in I}\E e^{\ii sJ_i+\ii vZ_i}\right|
          \le\exp\bigl(K_v|s|\sqrt A-cAs^2\bigr)\qquad(|s|\le c_0),
\]
whereas on $c_0\le|s|\le\pi$ the uniformly binary-like twisted core product is
exponentially small in $M$. Substituting $s=w/\sqrt A$, dominated convergence
and tightness from \eqref{eq:all-layer} together with \eqref{eq:integer-local}
yield \eqref{eq:joint-local}.
\end{proof}

\begin{lemma}[Asymptotic threshold envelopes]\label{lem:endpoints}
If $y/\sqrt A\to x\in\R$, let $j=\lfloor y+\mu\rfloor$ and suppose along a
subsequence that $\theta=y+\mu-j\to\theta_0\in[0,1]$ and $T\Rightarrow U$.
Then
\begin{align}
 \limsup_{n\to\infty}\sqrt A\,\bigl\{F_S(y)-\Phi(y/\sqrt V)\bigr\}
 &\le\phi(x)\left\{\frac12+\frac{1-2p}{6}(1-x^2)
                         -\E\{\theta_0-U\}\right\},\label{eq:closed-envelope}\\
 \liminf_{n\to\infty}\sqrt A\,\bigl\{F_S^-(y)-\Phi(y/\sqrt V)\bigr\}
 &\ge\phi(x)\left\{-\frac12+\frac{1-2p}{6}(1-x^2)\right\},
                                                        \label{eq:strict-envelope}
\end{align}
where $\{w\}=w-\lfloor w\rfloor$ denotes fractional part. If
$|y|/\sqrt A\to\infty$, then $\sqrt A\,|F_S^\pm(y)-\Phi(y/\sqrt V)|\to0$.
\end{lemma}
\begin{proof}
Decomposing $F_S(y)-F_K(j)=\sum_{\ell\in\Z}\{\Pp(K=j+\ell,T\le\theta-\ell)-\ind_{\{\ell\le0\}}\Pp(K=j+\ell)\}$,
truncating at $|\ell|\le H$ via \eqref{eq:all-layer} (which bounds the tail
by $KQ/(H\sqrt A)$), applying \eqref{eq:joint-local} and Portmanteau on closed
half-lines, and sending $H\to\infty$ gives
$\phi(x)\E\lfloor\theta_0-U\rfloor=\phi(x)(\theta_0-\E\{\theta_0-U\})$ (since
$\E U=0$). Combining this with \eqref{eq:integer-edgeworth} at $j=y+\mu-\theta$
proves \eqref{eq:closed-envelope}, and open half-lines similarly prove
\eqref{eq:strict-envelope}.
\end{proof}

\subsection{Modular concentration and proof of Theorem~\ref{thm:near}}

For a sequence satisfying \eqref{eq:near-assumption} with $p=\pe$, we have
$P/V\to\re$ and $\kappa/V\to\te$ by Proposition~\ref{prop:core}. The negative
strict branch \eqref{eq:strict-envelope} has normalized upper limit
$\phz(1/2-\te/6)/\re<\ce$, while the positive closed branch
\eqref{eq:closed-envelope} has unique maximum $\ce$ at $x=0$ and
$\E\{\theta_0-U\}=0$. Thus $U$ is concentrated on the integer coset
$\theta_0+\Z$, so there exist phases $\vartheta_n$ such that
\begin{equation}\label{eq:modular}
              \E\dist(S-\vartheta_n,\Z)\longrightarrow0.
\end{equation}

To deduce $\mathcal B(X_n)\to1$ from \eqref{eq:modular}, let $q=\qe$,
$t=2q-1=\te$, and let $P_q(u)$ be the $1$-periodic sawtooth function equal to
$-2(1-q)u$ on $[0,q]$ and $2q(u-1)$ on $[q,1]$. Then $P_q(x-\vartheta)$ has
mean $-q(1-q)$ on the circle and vanishes on $\vartheta+\Z$ (with
$|P_q(x-\vartheta)|\le2\dist(x-\vartheta,\Z)$), while the combined test
function $-tx+P_q(x-\vartheta)$ is $1$-Lipschitz on $\R$. Using
$\E S^*=\kappa/(2V)$ and $|\E e^{\ii uS^*}|\le1/(|u|\sqrt V)$ in the Fourier
series of $P_q(x-\vartheta)$ gives, by Kantorovich--Rubinstein duality,
\begin{equation}\label{eq:modular-dual}
 W_1(S,S^*)\ge\frac{t\kappa}{2V}+q(1-q)
  -2\E\dist(S-\vartheta,\Z)-\frac{\zeta(3)}{\pi^3\sqrt V}.
\end{equation}
Taking $\liminf$ via \eqref{eq:modular} gives
$\liminf W_1(S,S^*)\ge\te^2/2+\pe\qe=\re/2$, which matches the upper bound
$W_1(S,S^*)\le P/(2V)\to\re/2$ and proves $\mathcal B(X_n)\to1$.

Finally, \eqref{eq:modular} implies $|\E e^{2\pi\ii S}|=\prod_i|\E e^{2\pi\ii Z_i}|\to1$,
so $Q=\sum_i\E Z_i^2\le\frac14\sum_i(1-|\E e^{2\pi\ii Z_i}|^2)\to0$. Since
$|Z_i|\le1$, the direct couplings $(Y_i,J_i)$ satisfy
$\sum_i W_2^2(Y_i,J_i)\le Q\to0$ and $\sum_i W_3^3(Y_i,J_i)\le Q\to0$.
Restoring the gap scale $h_n$ proves \eqref{eq:strong-lattice} and completes
the proof of Theorem~\ref{thm:near}.

\begin{remark}
While $\mathcal B(X_n)\to1$ whenever $L_n\to0$ and $C_n\to\ce$, the pointwise
inequality $C\le\ce\mathcal B(X)$ does not hold for general distributions.
Indeed, for $N=1$, replacing the standard normal law $Z$ on $[-a,a]$ by three
atoms at $\{-a,0,a\}$ with matching mean and second moment (weights
$Pr/2,P(1-r),Pr/2$ where $P=\Pp(|Z|\le a)$ and $r=\E[Z^2\ind_{\{|Z|\le a\}}]/(a^2P)\to1/3$)
gives $d_{\mathrm K}=P(1-r)/2$ and $W_1=(aP/2)\{(1-r)^2+r^2\}$, so
$C/\mathcal B(X)=(1-r)/[2a\{(1-r)^2+r^2\}]\sim3/(5a)\to\infty$ as $a\downarrow0$.
\end{remark}
